\clearpage
\chapter*{第195次报告勘误}
\addcontentsline{toc}{chapter}{第195次报告勘误}
\sid{e195}
\setcounter{footnote}{0}
\src{300}

\noindent\textbf{4. TDTE II：形式模理论中的存在性定理。}

\begin{center}
[布尔巴基讨论班，第12卷，1959/60年度，第195号，22页。]
\end{center}

\noindent\textbf{第8页。} --- 命题5.1中所写的公式只有在 $\Lambda$ 是域时
才正确；一般情形下，必须把 $\mathfrak m_\xi/\mathfrak m_\xi^2$ 换成这个
空间对 $\mathfrak n_\xi/\mathfrak n_\xi^2$ 的像所取的商，其中
$\mathfrak n$ 是 $\Lambda$ 的极大理想。此外，关于 $O$ 在 $\Lambda$ 上
光滑所给出的定义，只有在剩余域扩张 $k'/k$ 可分时才正确。一般情形参见
SGA III，1.1。

\noindent\textbf{第14页，评注。} --- 这里提出的问题在射影情形下已由
“Hilbert 概形”彻底解决（参见 TDTE IV）。Nagata 和 Hironaka 的例子反而
表明，若不作射影性假设，所考察的函子未必可表；即便只限于分类一个三维完备
非奇异簇的零维子簇，情形也是如此。这与这类簇的对称平方可能不存在这一事实
有关。

\src{301}
\noindent\textbf{第15页，第3项。} --- 更完整的研究见 TDTE V。

\noindent\textbf{第16页，命题3.1。} --- 把“若 $f$ 固有”改为
“若 $f$ 固有且可分”。

\noindent\textbf{第18、19页，评注，第2项。} --- 我最近证明，域上
阿贝尔簇的形式模概形确实在 Witt 环上光滑；换言之，若一个阿廷局部环是
另一环的商，则该阿廷局部环上的任一阿贝尔概形，都是由后一环上的某个阿贝尔
概形约化而来。证明只用到第182次报告第12页所引入的提升障碍类的变换性质。
还应回顾，亏格为 $g$ 的曲线以及极化阿贝尔概形的模概形已由 Mumford 构造
（参见 SHMT）。

\noindent\textbf{第19页，第5段。} --- 为定义 $U$ 所考察的截面，必须假设
它不仅在 $\mathfrak X$ 上，而且在每个 $X_n$ 上都是非零因子。

\noindent\textbf{第20页，第12、13行。} --- 把
\[
  G_x=G\bigl(\operatorname{Spec}(\widehat{\mathcal O}_{\mathfrak X,x})\bigr),
\]
改为
\[
  G_x=G\bigl(\operatorname{Spec}(\mathcal O_{\mathfrak X,x})\bigr)
  \subset
  G\bigl(\operatorname{Spec}(\widehat{\mathcal O}_{\mathfrak X,x})\bigr),
\]
并把“$\mathcal O'_{0x}$”改为“$\widehat{\mathcal O}'_{0x}$”。
